\chapter{순전파, 역전파, optimizer를 한 흐름으로 보기}
\label{bg:chapter-forward-backward}

\section{순전파는 값을 만들고, 역전파는 책임을 나눈다}

순전파는 layer를 따라 activation을 계산한다. 두 파라미터 모델
$f_{w,b}(x)=wx+b$에서 $x=2$, $y=10$, $w=3$, $b=1$이면 $\hat y=7$이다.
제곱손실은 $\ell=\frac12(7-10)^2=4.5$다. 출력 오차 $-3$이 $w$와 $b$에
각각 얼마나 책임이 있는지 계산하는 과정이 \concept{backpropagation}{역전파}다.

\concept{gradient}{기울기}는 손실의 국소 경사다.

\begin{align}
  \frac{\partial \ell}{\partial w} &= (wx+b-y)x = (-3)\cdot2=-6,\\
  \frac{\partial \ell}{\partial b} &= (wx+b-y)=-3.
  \label{eq:bg-two-parameter-gradient}
\end{align}

음의 기울기는 해당 파라미터를 늘리면 손실이 줄어드는 국소 방향임을 뜻한다. 학습률
$\eta=0.1$인 가장 단순한 gradient descent라면

\begin{equation}
  w' = w-\eta(-6)=3.6,\qquad b'=b-\eta(-3)=1.3
\end{equation}

이고 새 예측은 $8.5$, 손실은 $1.125$로 줄어든다. 이 예는 ``기울기가 기억 그 자체''가
아님을 보여 준다. 기울기는 지금 본 예에 대한 국소 변경 제안이고, 실제로 어떤 상태를
얼마나 바꿀지는 optimizer가 정한다.

\section{Optimizer는 update rule과 기억 상태를 가진다}

\concept{optimizer}{옵티마이저}는 기울기를 파라미터 변화로 바꾸는 알고리즘이다.
SGD는 현재 기울기를 직접 쓰지만, momentum은 과거 기울기의 이동평균을 유지한다.
Adam 계열은 보통 1차 모멘트 $m_t$와 2차 모멘트 $v_t$를 유지한다.
\concept{learning-rate}{학습률} $\eta$는 한 update에서 이 방향을 얼마나 크게 반영할지
정하는 scale이다.

\begin{align}
  m_t &= \beta_1m_{t-1}+(1-\beta_1)g_t,\\
  v_t &= \beta_2v_{t-1}+(1-\beta_2)g_t^2,\\
  \theta_{t+1} &= \theta_t-\eta
  \frac{\hat m_t}{\sqrt{\hat v_t}+\epsilon}.
  \label{eq:bg-adam}
\end{align}

$m_t$와 $v_t$가 \concept{optimizer-state}{옵티마이저 상태}다. parameter가 BF16으로
2 byte라 해도, 안정적인 학습을 위해 FP32 master weight, gradient, 두 moment를 보관하면
학습 중 persistent state는 parameter당 12--16 byte 수준이 될 수 있다. sharding과
8-bit optimizer가 이를 줄일 수 있지만, ``LoRA parameter가 작다''와 ``sleep job 전체
memory footprint가 작다''는 동일한 문장이 아니다. activation과 temporary buffer가 별도로
존재하기 때문이다.

\begin{table}[htbp]
\centering
\caption{추론과 학습 단계의 대표 상태}
\label{tab:bg-training-state}
\begin{tabularx}{\textwidth}{L{0.21\textwidth}YYY}
\toprule
상태 & 추론 & 학습 & 수명 \\
\midrule
model weights & 읽기 중심 & 읽기+갱신 & checkpoint 간 지속 \\
activations & 현재 token/layer & backward까지 보존 또는 재계산 & step 내부 \\
gradients & 없음 & trainable parameter마다 생성 & accumulation window \\
optimizer state & 없음 & momentum/variance/master copy & 여러 step 지속 \\
data state & prompt/KV & shuffle cursor, sampler, replay provenance & run 전체 \\
validation artifact & 선택적 & metric, failure slice, approval & release와 함께 지속 \\
\bottomrule
\end{tabularx}
\end{table}

\section{Activation memory와 checkpointing}

역전파는 각 layer의 입력 activation이 있어야 기울기를 계산하기 쉽다. 모든 activation을
HBM에 보관하면 빠르지만 memory를 많이 쓴다. 일부를 버리고 backward 때 다시 순전파하는
activation checkpointing은 compute를 늘려 memory를 줄인다. 따라서 같은 training FLOPs라
해도 sequence length, microbatch, layer checkpoint policy에 따라 HBM 요구량과 wall-clock이
달라진다.

Sleep-time training은 wake serving과 같은 GPU에서 돌아갈 수 있으므로 peak memory만이
문제가 아니다. checkpoint save가 storage bandwidth를 점유하고, parameter all-gather가
interconnect를 점유하며, validation이 inference slot을 가져간다. 좋은 스케줄러는 sleep job의
평균 utilization이 아니라 wake p99 latency를 보호해야 한다.

\section{Gradient accumulation과 effective batch}

한 번에 큰 batch가 HBM에 들어가지 않으면 여러 microbatch의 gradient를 더한 뒤 한 번
update한다. accumulation step을 $K$, microbatch 크기를 $B_\mu$라 하면 data-parallel replica
하나의 local effective batch는 $K B_\mu$다. replica가 $N$개이고 각 replica의 gradient를
동기화하면 global batch는 대략 $NKB_\mu$다.

하지만 token 길이가 다르면 ``문장 수 batch''와 ``token 수 batch''는 다르다. memory
consolidation에서는 긴 episode 몇 개가 gradient와 activation 대부분을 만들 수 있다.
따라서 scheduler는 examples/s보다 tokens/s, effective valid memories/s, update당 unique
provenance를 함께 기록해야 한다.

\begin{cautionbox}[작은 loss가 곧 좋은 장기기억은 아니다]
Training loss가 내려가도 replay 샘플의 중복, contamination, teacher 오류를 외웠을 수 있다.
반대로 old-task loss를 너무 강하게 고정하면 새 정보가 들어갈 공간이 사라진다. backward는
정확히 계산할 수 있지만, 무엇을 backward할지 정하는 목적과 데이터가 더 근본적인 문제다.
\end{cautionbox}

\section{Sleep-time operator로 번역하기}

Parametric sleep job을 한 번 실행한다고 하자. 먼저 selected episodes를 forward해 loss를
계산하고, backward로 gradient를 만든 뒤 optimizer가 adapter 또는 weights를 갱신한다.
그 다음 held-out memory query, old-capability suite, privacy probe, latency benchmark를
실행한다. 새 artifact가 모두 통과하면 version pointer를 바꾸고, 실패하면 이전 pointer를
유지한다. 이 순서에서 학습 kernel보다 중요한 시스템 계약은 다음과 같다.

\begin{enumerate}
  \item source episode와 gradient-producing example 사이의 provenance가 유지될 것,
  \item training state가 wake serving state와 원자적으로 분리될 것,
  \item validation 전 artifact가 사용자 traffic에 노출되지 않을 것,
  \item optimizer state를 재개할지 폐기할지 명시할 것,
  \item 삭제 요청 시 derived checkpoint와 adapter까지 추적할 수 있을 것.
\end{enumerate}
